\chapter{Introduzione}

\section{Ripasso Algebra Lineare}
\begin{defn}[Spazio vettoriale]
  Lo \textbf{spazio vettoriale} $\mathbb{R}^n$ è un insieme di elementi $x = [x_1, \cdots, x_n]$ where each $x_i \in \mathbb{R}$.  
  Gli elementi $x$ sono chiamati \textbf{vettori} e soddisfano queste proprietà:
  \begin{enumerate}
    \item \textbf{Addizione:} Se $x \in \mathbb{R}^n$ e $y \in \mathbb{R}^n$, allora $x + y = [x_1+y_1, \cdots, x_n+y_n] \in \mathbb{R}^n$. 
    \item \textbf{Prodotto scalare:} Se $x \in \mathbb{R}^n$ e $a \in \mathbb{R}$ allora, $a \cdot x = [a\cdot x_1, \cdots, a\cdot x_n]$.
    \item \textbf{Prodotto interno:} Se $x \in \mathbb{R}^n$ e $y \in \mathbb{R}^n$ allora, $x \cdot y = \sum^n_{i=1}{x_i \cdot y_i}$.
  \end{enumerate}
\end{defn}

\begin{defn}[Matrice]
  Una \textbf{matrice} $\text{X} \in \mathbb{R}^{n \times m}$ è un vettore rettangolare di elementi $x_{ij} \in \mathbb{R}$, $1 \le i \le n$, $1 \le j \le m$:

  \[\mathbf{X} = \begin{pmatrix}
    x_{11} & x_{12} & \cdots & x_{1m} \\
    x_{21} & x_{22} & \cdots & x_{2m} \\
    \vdots & \vdots & \ddots & \vdots \\
    x_{n1} & x_{n2} & \cdots & x_{nm}
  \end{pmatrix}\]

  Deve supportare queste operazioni:
  \begin{enumerate}
    \item \textbf{Addizione}: Se $\text{X} \in \mathbb{R}^{n \times m}$, $\text{Y} \in \mathbb{R}^{n \times m}$ e $\text{Z} = \text{X} + \text{Y}$, allora $\text{Z} \in \mathbb{R}^{n \times m}$ e $Z_{ij} = X_{ij} + Y_{ij}$.
    \item \textbf{Prodotto scalare}: Se $\text{X} \in \mathbb{R}^{n \times m}$, $a \in \mathbb{R}$ e $\text{Z} = a \cdot \text{X}$, allora $\text{Z} \in \mathbb{R}^{n \times m}$ e $Z_{ij} = a \cdot X_{ij}$.
    \item \textbf{Prodotto fra matrici}: Se $\text{X} \in \mathbb{R}^{n \times m}$, $\text{Y} \in \mathbb{R}^{m \times n}$ e $\text{Z} = \text{X} \cdot \text{Y}$, allora $\text{Z} \in \mathbb{R}^{n \times n}$ e $Z_{ij} = \sum_{k=1}^{m} X_{ik} \cdot Y_{kj}$.
  \end{enumerate}
\end{defn}

\section{Ripasso Statistica}
\begin{defn}[Distribuzione di Probabilità]
  Una distribuzione di probabilità $P$ su uno spazio campionario $\Omega$ è un'applicazione dai sottoinsiemi di $\Omega$ ai numeri reali che soddisfa le seguenti condizioni:
  \begin{itemize}
    \item \textbf{Non-negatività:} $P(\alpha) \ge 0$, $\forall \alpha \subseteq \Omega$.
    \item \textbf{Normalizzazione:} $P(\Omega) = 1$.
    \item \textbf{Additività:} $\forall \alpha, \beta \subseteq \mathbb{R}$ che sono insiemi disgiunti, $P(\alpha \cup \beta) = P(\alpha) + P(\beta)$.
  \end{itemize}
\end{defn}


\section{Machine Learning}
\begin{defn}[Machine Learning]
Il \textbf{machine learning} è un campo di studio
che conferisce ai computer la capacità di apprendere senza essere
esplicitamente programmati.
\end{defn}

Esistono due tipi di approcci:
\begin{enumerate}
  \item \textbf{Non parametrico:} l'apprendimento avviene memorizzando i dati di addestramento (memorizzazione).
  \item \textbf{Parametrico:} l'apprendimento avviene utilizzando un algoritmo per adattare i parametri di un modello matematico o statistico sulla base di dati di addestramento. Per esempio:
    \[
    f_{\theta}(x) = \sum^D_{d=1}{\theta_d x_d}
    \]
\end{enumerate}


